\documentclass[10pt]{extarticle}
\usepackage[a4paper,margin=1.5cm]{geometry}
\usepackage[T1]{fontenc}
\usepackage{amsmath,amssymb}
\usepackage{graphicx}
\usepackage[export]{adjustbox}
\usepackage{multicol}
\usepackage{enumitem}
\usepackage{booktabs,array,tabularx}
\setlength{\columnsep}{0.7cm}
\setlength{\parindent}{0pt}
\setlength{\parskip}{3pt}
\pagestyle{empty}
\makeatletter
% figure registry (figdefs.tex is generated by the extract step)
\newcommand{\qfigdef}[4]{\@namedef{qf@file@#1}{#2}\@namedef{qf@w@#1}{#3}\@namedef{qf@wide@#1}{#4}}
\newcommand{\qf@check}[1]{\@ifundefined{qf@file@#1}{\PackageError{score2}{Unknown figure id #1}{}}{}}
\newcommand{\qf@img}[4][]{\includegraphics[width=#2\dimexpr\@nameuse{qf@w@#4}pt\relax,max width=#3,max height=0.45\textheight,#1]{\@nameuse{qf@file@#4}}}
% Full-width questions: finalize wraps every question that contains a wide figure in \qfullbegin...\qfullend,
% so the WHOLE question (stem, figures, options) is set across the page in reading order, then the two columns
% resume. The question's head - its number, stem and first figure (\qfullsplit marks the end of the head) -
% is kept on one page, so a number is never left at a page foot with its figure overleaf; the rest may break
% across pages normally, so no big gaps are left. Without a \qfullsplit, the whole question is kept together
% if it is under 45% of a page. Never write these yourself.
\newif\ifqf@full
\newif\ifqf@boxed
\newsavebox{\qf@box}
\newcommand{\qf@open}{\global\qf@fulltrue\par\global\qf@boxedtrue
  \setbox\qf@box\vbox\bgroup\hsize\textwidth\linewidth\textwidth}
\newcommand{\qf@emit}[1]{\par\egroup\global\qf@boxedfalse
  \ifdim\dimexpr\ht\qf@box+\dp\qf@box\relax<#1\textheight\noindent\usebox{\qf@box}\par\else\unvbox\qf@box\fi}
\newcommand{\qfullsplit}{\ifqf@boxed\qf@emit{0.85}\fi}
\newcommand{\qf@close}{\ifqf@boxed\qf@emit{0.45}\fi\global\qf@fullfalse}
\newcommand{\qfullbegin}{\end{multicols}\qf@open}
\newcommand{\qfullend}{\qf@close\begin{multicols}{2}}
\newcommand{\qfullstart}{\qf@open}
\newcommand{\qfulldone}{\qf@close}
% \qfig[scale]{id}: block figure, centred, directly where it is written (size set by the tool).
% Fallback if a wide figure is not inside a full-width question: break out of the columns in place - no floats.
\newcommand{\qfig}[2][1]{\qf@check{#2}%
  \ifqf@full\qf@blk{#1}{#2}\else
  \ifnum\@nameuse{qf@wide@#2}=1 \expandafter\@firstoftwo\else\expandafter\@secondoftwo\fi
  {\end{multicols}\begin{center}\qf@img{#1}{\textwidth}{#2}\end{center}\begin{multicols}{2}}%
  {\qf@blk{#1}{#2}}\fi}
\newcommand{\qf@blk}[2]{\par\smallskip{\centering\qf@img{#1}{\linewidth}{#2}\par}\smallskip}
% \qfigi[scale]{id}: inline figure (inside option lists / tables / minipages); never breaks columns.
\newcommand{\qfigi}[2][1]{\qf@check{#2}\qf@img[valign=c]{#1}{\linewidth}{#2}}
\makeatother
% structure macros - lint reads these to check completeness and figure placement
\newcommand{\papersec}[1]{\par\bigskip{\centering\Large\bfseries PAPER-#1\par}\medskip}
\newcommand{\partsec}[2]{\par\medskip\noindent{\large\bfseries PART-#1: #2}\par\smallskip}
\newcommand{\sect}[1]{\par\smallskip\noindent\textbf{#1}\par}
\newcommand{\qn}[1]{\par\smallskip\noindent\textbf{#1.}~}
\newcommand{\paraq}[2]{\par\smallskip\noindent\textbf{Paragraph for Questions #1 and #2}\par}
\newenvironment{qoptions}{\begin{itemize}[leftmargin=1.8em,itemsep=1pt,topsep=2pt,label={}]}{\end{itemize}}
\newcommand{\opt}[2]{\item[(#1)] #2}
% figure registry (generated by score2_tools.py finalize)
\qfigdef{p1_phy_s1_q04_1}{Nurture (General) - Test 03 (27 Sep 2023)-img/p1_phy_s1_q04_1.png}{104.1}{0}
\qfigdef{p1_phy_s1_q05_A}{Nurture (General) - Test 03 (27 Sep 2023)-img/p1_phy_s1_q05_A.png}{78.6}{0}
\qfigdef{p1_phy_s1_q05_B}{Nurture (General) - Test 03 (27 Sep 2023)-img/p1_phy_s1_q05_B.png}{77.6}{0}
\qfigdef{p1_phy_s1_q05_C}{Nurture (General) - Test 03 (27 Sep 2023)-img/p1_phy_s1_q05_C.png}{74.8}{0}
\qfigdef{p1_phy_s1_q05_D}{Nurture (General) - Test 03 (27 Sep 2023)-img/p1_phy_s1_q05_D.png}{75.8}{0}
\qfigdef{p1_phy_s1_q08_1}{Nurture (General) - Test 03 (27 Sep 2023)-img/p1_phy_s1_q08_1.png}{150.2}{0}
\qfigdef{p1_phy_s1_q11_A}{Nurture (General) - Test 03 (27 Sep 2023)-img/p1_phy_s1_q11_A.png}{54.4}{0}
\qfigdef{p1_phy_s1_q11_B}{Nurture (General) - Test 03 (27 Sep 2023)-img/p1_phy_s1_q11_B.png}{54.6}{0}
\qfigdef{p1_phy_s1_q11_C}{Nurture (General) - Test 03 (27 Sep 2023)-img/p1_phy_s1_q11_C.png}{54.1}{0}
\qfigdef{p1_phy_s1_q11_D}{Nurture (General) - Test 03 (27 Sep 2023)-img/p1_phy_s1_q11_D.png}{54.1}{0}
\qfigdef{p1_phy_s1_q15_1}{Nurture (General) - Test 03 (27 Sep 2023)-img/p1_phy_s1_q15_1.png}{92.1}{0}
\qfigdef{p1_phy_s1_q19_1}{Nurture (General) - Test 03 (27 Sep 2023)-img/p1_phy_s1_q19_1.png}{115.6}{0}
\qfigdef{p1_phy_s2_q03_1}{Nurture (General) - Test 03 (27 Sep 2023)-img/p1_phy_s2_q03_1.png}{114.4}{0}
\qfigdef{p1_phy_s2_q05_1}{Nurture (General) - Test 03 (27 Sep 2023)-img/p1_phy_s2_q05_1.png}{117.1}{0}
\qfigdef{p1_chem_s1_q05_1}{Nurture (General) - Test 03 (27 Sep 2023)-img/p1_chem_s1_q05_1.png}{103.5}{0}
\qfigdef{p1_chem_s1_q08_A}{Nurture (General) - Test 03 (27 Sep 2023)-img/p1_chem_s1_q08_A.png}{85.9}{0}
\qfigdef{p1_chem_s1_q08_B}{Nurture (General) - Test 03 (27 Sep 2023)-img/p1_chem_s1_q08_B.png}{84.1}{0}
\qfigdef{p1_chem_s1_q08_C}{Nurture (General) - Test 03 (27 Sep 2023)-img/p1_chem_s1_q08_C.png}{86.9}{0}
\qfigdef{p1_chem_s1_q08_D}{Nurture (General) - Test 03 (27 Sep 2023)-img/p1_chem_s1_q08_D.png}{81.9}{0}
\qfigdef{p1_chem_s1_q14_1}{Nurture (General) - Test 03 (27 Sep 2023)-img/p1_chem_s1_q14_1.png}{78.1}{0}
\qfigdef{p1_chem_s1_q20_1}{Nurture (General) - Test 03 (27 Sep 2023)-img/p1_chem_s1_q20_1.png}{107.5}{0}
\qfigdef{p1_chem_s2_q04_1}{Nurture (General) - Test 03 (27 Sep 2023)-img/p1_chem_s2_q04_1.png}{141.4}{0}
\raggedcolumns
\makeatletter\@beginparpenalty=10000\makeatother
\newcolumntype{L}[1]{>{\raggedright\arraybackslash}p{#1}}
\begin{document}
\begin{center}{\Large\bfseries Nurture (General) --- Test 03 (27 Sep 2023)}\end{center}
\begin{multicols}{2}
\raggedcolumns
\papersec{1}
\partsec{1}{PHYSICS}
\sect{SECTION-I}
\qn{1} An inductance $L$ and a resistance $R$ are first connected to a battery. After some time (assume this time as $t=0$) the battery is disconnected but $L$ and $R$ remain connected in a closed circuit. Then the current reduces to 37\% of its initial value in
\begin{qoptions}
\opt{A}{$RL$ sec}
\opt{B}{$\dfrac RL$ sec}
\opt{C}{$\dfrac LR$ sec}
\opt{D}{$\dfrac1{LR}$ sec}
\end{qoptions}
\qn{2} A coil of inductance 300 $mH$ and resistance 2$\Omega$ is connected to a source of voltage 2V. The current reaches half of its steady state value in
\begin{qoptions}
\opt{A}{0.15 s}
\opt{B}{0.3 s}
\opt{C}{0.05 s}
\opt{D}{0.1 s}
\end{qoptions}
\qn{3} An e.m.f. of 15 volt is applied in a circuit containing 5 henry inductance and 10 ohm resistance. The ratio of the currents at time $t=\infty$ and at $t=1$ second is
\begin{qoptions}
\opt{A}{$\dfrac{e^{1/2}}{e^{1/2}-1}$}
\opt{B}{$\dfrac{e^2}{e^2-1}$}
\opt{C}{$1-e^{-1}$}
\opt{D}{$e^{-1}$}
\end{qoptions}
\qn{4} When a certain circuit consisting of a constant e.m.f. $E$ an inductance $L$ and a resistance $R$ is closed, the current in, it increases with time according to curve 1. After one parameter ($E$, $L$ or $R$) is changed, the change in current follows curve 2 when the circuit is closed second time. Which parameter was changed and in what direction\par
\qfig{p1_phy_s1_q04_1}
\begin{qoptions}
\opt{A}{$L$ is increased}
\opt{B}{$L$ is decreased}
\opt{C}{$R$ is increased}
\opt{D}{$R$ is decreased}
\end{qoptions}
\qn{5} In an L--R circuit connected to a battery the rate at which energy is stored in the inductor is plotted against time during the growth of the current in the circuit. Which of the following best represents the resulting curve
\begin{qoptions}
\opt{A}{\qfigi{p1_phy_s1_q05_A}}
\opt{B}{\qfigi{p1_phy_s1_q05_B}}
\opt{C}{\qfigi{p1_phy_s1_q05_C}}
\opt{D}{\qfigi{p1_phy_s1_q05_D}}
\end{qoptions}
\qn{6} A coil of 160 turns of cross - sectional area 250 cm$^2$ rotates at an angular velocity of 300 rad/sec about an axis parallel to the plane of the coil in a uniform magnetic field of 0.6 weber/meter$^2$. If the coil is connected to a resistance of 2 ohm what is the maximum torque that has to be delivered to maintain its motion ?
\begin{qoptions}
\opt{A}{650 N $-$ m}
\opt{B}{864 N $-$ m}
\opt{C}{720 N $-$ m}
\opt{D}{360 N $-$ m}
\end{qoptions}
\qn{7} An inductor $L$ is allowed to discharge through a capacitor $C$. Find the emf induced across the inductor when the capacitor is fully charged.
\begin{qoptions}
\opt{A}{Maximum}
\opt{B}{Minimum}
\opt{C}{Zero}
\opt{D}{Infinite}
\end{qoptions}
\qn{8} The voltage time (V-t) graph for triangular wave having peak value $V_0$ is as shown in figure. The rms value of $V$ in time interval from $t=0$ to $T/4$ is $\dfrac{V_0}{\sqrt x}$ then find the value of $x$.\par
\qfig{p1_phy_s1_q08_1}
\begin{qoptions}
\opt{A}{5}
\opt{B}{4}
\opt{C}{7}
\opt{D}{3}
\end{qoptions}
\qn{9} The instantaneous values of alternating current and voltages in a circuit are given as\par
$i=\dfrac1{\sqrt2}\sin(100\pi t)$ampere,\par
$e=\dfrac1{\sqrt2}\sin(100\pi t+\pi/3)$volt\par
The average power in Watts consumed in the circuit is :
\begin{qoptions}
\opt{A}{$\dfrac14$}
\opt{B}{$\dfrac{\sqrt3}4$}
\opt{C}{$\dfrac12$}
\opt{D}{$\dfrac18$}
\end{qoptions}
\qn{10} In an alternating current circuit in which an inductance and capacitance are joined in series, current is found to be maximum when the value of inductance is 0.5 henry and the value of capacitance is 8$\mu$F. The angular frequency of applied alternating voltage will be
\begin{qoptions}
\opt{A}{5000 rad/sec}
\opt{B}{4000 rad/sec}
\opt{C}{$2\times10^5$ rad/sec}
\opt{D}{500 rad/sec}
\end{qoptions}
\qn{11} Which one of the following curves represents the variation of impedance ($Z$) with frequency $f$ in series LCR circuit?
\begin{qoptions}
\opt{A}{\qfigi{p1_phy_s1_q11_A}}
\opt{B}{\qfigi{p1_phy_s1_q11_B}}
\opt{C}{\qfigi{p1_phy_s1_q11_C}}
\opt{D}{\qfigi{p1_phy_s1_q11_D}}
\end{qoptions}
\qn{12} A coil of resistance 50$\Omega$ is connected across a 5.0 V battery. 0.1 s after the battery is connected, the current in the coil is 60 mA. Calculate the inductance of the coil,
\begin{qoptions}
\opt{A}{5.5 H}
\opt{B}{1.5 H}
\opt{C}{2.5 H}
\opt{D}{9.5 H}
\end{qoptions}
\qn{13} An LCR series circuit is connected to a source of alternating current. At resonance, the applied voltage and the current flowing through the circuit will have a phase difference of
\begin{qoptions}
\opt{A}{$\pi$}
\opt{B}{$\dfrac\pi2$}
\opt{C}{$\dfrac\pi4$}
\opt{D}{0}
\end{qoptions}
\qn{14} In an electrical circuit R, L, C and an a.c. voltage source are all connected in series. When L is removed from the circuit, the phase difference between the voltage the current in the circuit is $\pi/3$. If instead, C is removed from the circuit, the phase difference is again $\pi/3$. The power factor of the circuit is
\begin{qoptions}
\opt{A}{$\dfrac12$}
\opt{B}{$\dfrac1{\sqrt2}$}
\opt{C}{1}
\opt{D}{$\dfrac{\sqrt3}2$}
\end{qoptions}
\qn{15} In an LCR circuit as shown below both switches are open initially. Now switch $S_1$ is closed, $S_2$ kept open, ($q$ is charge on the capacitor and $\tau=RC$ is Capacitive time constant). Which of the following statement is correct?\par
\qfig{p1_phy_s1_q15_1}
\begin{qoptions}
\opt{A}{Work done by the battery is half of the energy dissipated in the resistor}
\opt{B}{At $t=\tau$, $q=CV/2$}
\opt{C}{At $t=2\tau$, $q=CV\left(1-e^{-2}\right)$}
\opt{D}{At $t=2\tau$, $q=CV\left(1-e^{-1}\right)$}
\end{qoptions}
\qn{16} If $\vec E$ and $\vec B$ represent electric and magnetic field vectors of the electromagnetic wave, the direction of propagation of electromagnetic wave is along
\begin{qoptions}
\opt{A}{$\vec E$}
\opt{B}{$\vec B$}
\opt{C}{$\vec B\times\vec E$}
\opt{D}{$\vec E\times\vec B$}
\end{qoptions}
\qn{17} A plane electromagnetic wave travels in free space along x-axis. At a particular point in space, the electric field along y-axis is 9.3 V m$^{-1}$. The magnetic induction (B) along z-axis is
\begin{qoptions}
\opt{A}{$3.1\times10^{-8}$T}
\opt{B}{$3\times10^{-5}$T}
\opt{C}{$3\times10^{-6}$T}
\opt{D}{$9.3\times10^{-6}$T}
\end{qoptions}
\qn{18} The energy of electromagnetic wave in vacuum is given by the relation
\begin{qoptions}
\opt{A}{$\dfrac{E^2}{2\varepsilon_0}+\dfrac{B^2}{2\mu_0}$}
\opt{B}{$\dfrac12\varepsilon_0E^2+\dfrac12\mu_0B^2$}
\opt{C}{$\dfrac{E^2+B^2}{c}$}
\opt{D}{$\dfrac12\varepsilon_0E^2+\dfrac{B^2}{2\mu_0}$}
\end{qoptions}
\qn{19} A source of constant voltage $V$ is connected to a resistance $R$ and two ideal inductors $L_1$ and $L_2$ through a switch $S$ as shown. There is no mutual inductance between the two inductors. The switch $S$ is initially open. At $t=0$, the switch is closed and current begins to flow. Which of the following options is incorrect?\par
\qfig{p1_phy_s1_q19_1}
\begin{qoptions}
\opt{A}{At $t=0$, the current through the resistance $R$ is $\dfrac VR$}
\opt{B}{After a long time, the current through $L_2$ will be $\dfrac{V\,L_1}{R(L_1+L_2)}$}
\opt{C}{After a long time, the current through $L_1$ will be $\dfrac{V\,L_2}{R(L_1+L_2)}$}
\opt{D}{The ratio of the currents through $L_1$ and $L_2$ is fixed at all times ($t>0$).}
\end{qoptions}
\qn{20} The value of current in two series L C R circuits at resonance is same. Then
\begin{qoptions}
\opt{A}{Both circuits must have same value of capacitance and inductance}
\opt{B}{In both circuits ratio of L and C will be same}
\opt{C}{For both the circuits $X_L/X_C$ must be same at that frequency}
\opt{D}{Both circuits must have same impedance at all frequencies}
\end{qoptions}
\sect{SECTION-II}
\qn{1} A small town with a demand of 800 kW of electric power at 220 V is situated 15 km away from an electric plant generating power at 440 V. The resistance of the two wire line carrying power is 0.5 $\Omega$/km.. The town gets power from the line through a 4000-220 V step-down transformer at a sub-station in the town. How much power (In MW) must the plant supply, assuming there is negligible power loss due to leakage?
\qn{2} Radio receiver receives a message at 300m band, If the available inductance is 1 mH, then calculate required capacitance (in pF)\par
Radio receives EM waves. (velocity of EM waves $c=3\times10^8$m/s)
\qn{3} A horizontal copper disc of diameter 20 cm, makes 10 revolutions/sec about a vertical passing through its centre. A uniform magnetic field of 100 gauss acts perpendicular to the plane of the disc. The potential difference its centre and rim in volts is $n$(in mV) then find $n$.\par
\qfig{p1_phy_s2_q03_1}
\qn{4} A solenoid has 2000 turns wound over a length of 0.3 m. The area of cross-section is $1.2\times10^{-3}$m$^2$. Around its central section a coil of 300 turns is closely would. If an initial current of 2A is reversed in 0.25 s, find the emf induced in the coil. (In mV)
\qn{5} If $X_L=50\ \Omega$ and $X_C=40\ \Omega$. Calculate effective value of current in given circuit (in Amp) $\sqrt2=1.41$\par
\qfig{p1_phy_s2_q05_1}
\qn{6} When 10V, DC is applied across a coil current through it is 2.5 A, if 10V, 50 Hz A.C. is applied current reduces to 2 A. Calculate reactance of the coil (in ohm)
\qn{7} When an alternating voltage of 220V is applied across a device X, a current of 0.5 A flows through the circuit and is in phase with the applied voltage. When the same voltage is applied across another device Y, the same current again flows through the circuit but it leads the applied voltage by $\pi/2$ radians.\par
Calculate the current flowing in the circuit when same voltage is applied across the series combination of X and Y (in Amp)
\qn{8} An electric lamp which runs at 80V DC consumes 10 A current. The lamp is connected to 100 V $-$ 50 Hz ac source. Compute the inductance of the choke required. (in henry)
\qn{9} The magnetic field in a travelling electromagnetic wave has a peak value of 20 nT. The peak value of electric field strength is (in V/m)
\qn{10} A 27 mW laser beam has a cross-sectional area of 10 mm$^2$. The magnitude of the maximum electric field in this electromagnetic wave is given by [Given permittivity of space $\varepsilon_0=9\times10^{-12}$ SI units, speed of light $c=3\times10^8$m/s] in k V/m
\partsec{2}{CHEMISTRY}
\sect{SECTION-I}
\qn{1} Calculate the cell EMF in mV for\par
Pt $|$H$_2$(1 atm)$|$HC$\ell$(0.01 M)$|$AgC$\ell$(s)$|$Ag(s) at 298K\par
If $\Delta G^\circ_f$ values are at 25$^\circ$C.\par
$-109.56\,\dfrac{\text{kJ}}{\text{mol}}$ for AgC$\ell$(s) and $-130.79\,\dfrac{\text{kJ}}{\text{mol}}$ for (H$^+$ + C$\ell^-$)(aq)
\begin{qoptions}
\opt{A}{456 mV}
\opt{B}{654 mV}
\opt{C}{546 mV}
\opt{D}{None of these}
\end{qoptions}
\qn{2} Consider the cell Ag(s) $|$ AgBr(s) $|$ Br$^-$(aq) $||$ AgC$\ell$(s) $|$ C$\ell^-$(aq)$|$ Ag(s) at 25$^\circ$C. The solubility product constants of AgBr \& AgC$\ell$ are respectively $5\times10^{-13}$ \& $1\times10^{-10}$. For what ratio of the concentrations of Br$^-$ \& C$\ell^-$ ions would emf of the cell be zero?
\begin{qoptions}
\opt{A}{1:500}
\opt{B}{1:100}
\opt{C}{1:200}
\opt{D}{200:1}
\end{qoptions}
\qn{3} If $E^\circ_{\mathrm{Fe}^{+2}/\mathrm{Fe}}=-0.44\,$V and $E^\circ_{\mathrm{Fe}^{+3}/\mathrm{Fe}^{+2}}=0.77$. What is $E^\circ_{\mathrm{Fe}^{+3}/\mathrm{Fe}}$
\begin{qoptions}
\opt{A}{-0.036V}
\opt{B}{-0.33 V}
\opt{C}{+0.33V}
\opt{D}{+0.0367V}
\end{qoptions}
\qn{4} Electrolysis of dil. H$_2$SO$_4$ generates H$_2$S$_2$O$_8$ (per disulphuric acid). What current is needed to produce H$_2$S$_2$O$_8$ at the rate of 1 mole per hour assuming 75\% charge efficiency.
\begin{qoptions}
\opt{A}{53.7amp}
\opt{B}{17.48 amp}
\opt{C}{35.74 amp}
\opt{D}{71.48 amp}
\end{qoptions}
\qn{5} Two weak acid solution HA$_1$ and HA$_2$ each with the same concentration and having pK$_a$ values 3 and 5 are placed in contact with hydrogen electrodes (1atm,25$^\circ$C) and are interconnected through a salt bridge. The emf of the cell is: (Assuming $\alpha\ll1$ for both acids)\par
\qfig{p1_chem_s1_q05_1}
\begin{qoptions}
\opt{A}{0.21V}
\opt{B}{0.018V}
\opt{C}{0.059V}
\opt{D}{0.021 V}
\end{qoptions}
\qn{6} During the preparation of H$_2$S$_2$O$_8$ (per disulphuric acid) O$_2$ gas also releases at anode as by product, When 9.72 L of H$_2$ releases at cathode and 2.35 L O$_2$ at anode at STP, the weight of H$_2$S$_2$O$_8$ produced in gram is
\begin{qoptions}
\opt{A}{87.12}
\opt{B}{51.74}
\opt{C}{83.42}
\opt{D}{43.47}
\end{qoptions}
\qn{7} A galvanic cell is composed of two hydrogen electrodes, one of which is a standard one. In which of the following solutions should the other electrode be immersed to get maximum emf ? K$_a$(CH$_3$COOH) $=2\times10^{-5}$, K$_a$(H$_3$PO$_4$) $=10^{-3}$
\begin{qoptions}
\opt{A}{0.1 M HCl}
\opt{B}{0.1 M CH$_3$COOH}
\opt{C}{0.1 M H$_3$PO$_4$}
\opt{D}{0.1 M H$_2$SO$_4$}
\end{qoptions}
\qn{8} Conductance measurements can be used to detect the end point of acid-base titrations. Which of the following plots correctly represent the end point of the titration of strong acid and a strong base ?
\begin{qoptions}
\opt{A}{\qfigi{p1_chem_s1_q08_A}}
\opt{B}{\qfigi{p1_chem_s1_q08_B}}
\opt{C}{\qfigi{p1_chem_s1_q08_C}}
\opt{D}{\qfigi{p1_chem_s1_q08_D}}
\end{qoptions}
\qn{9} When a lead storage battery is discharged?
\begin{qoptions}
\opt{A}{PbSO$_4$ is formed}
\opt{B}{Pb is formed}
\opt{C}{SO$_2$ is consumed}
\opt{D}{H$_2$SO$_4$ is formed}
\end{qoptions}
\qn{10} The standard reduction potentials of some half cell reactions are given below :\par
PbO$_2$ + 4H$^+$ + 2e$^-\rightleftharpoons$ Pb$^{2+}$ + 2H$_2$O \quad $E^0=1.455$ V\par
MnO$_4^-$ + 8H$^+$ + 5e$^-\rightleftharpoons$ Mn$^{2+}$ + 4H$_2$O \quad $E^0=1.51$ V\par
Ce$^{4+}$ + e$^-\rightleftharpoons$ Ce$^{3+}$ \quad $E^0=1.61$ V\par
H$_2$O$_2$ + 2H$^+$ + 2e$^-\rightleftharpoons$ 2H$_2$O \quad $E^0=1.71$ V\par
Pick out the \textbf{incorrect} statement (for standard condition of all species):
\begin{qoptions}
\opt{A}{Ce$^{+4}$ will oxidise Pb$^{2+}$ to PbO$_2$}
\opt{B}{MnO$_4^-$ will oxidise Pb$^{2+}$ to PbO$_2$}
\opt{C}{H$_2$O$_2$ will oxidise Mn$^{+2}$ to MnO$_4^-$}
\opt{D}{PbO$_2$ will oxidise Mn$^{+2}$ to MnO$_4^-$}
\end{qoptions}
\qn{11} On electrolysis, in which of the following, O$_2$ would \textbf{not} be liberated at the anode ?
\begin{qoptions}
\opt{A}{dilute H$_2$SO$_4$ with Pt electrodes}
\opt{B}{aqueous AgNO$_3$ solution with Pt electrodes}
\opt{C}{dilute H$_2$SO$_4$ with Cu electrodes}
\opt{D}{aqueous NaOH with a Fe cathode \& a Pt anode}
\end{qoptions}
\qn{12} For the reaction 2NO (g) + H$_2$ (g) $\to$ N$_2$O (g) + H$_2$O (g) at 900K. following data are observed.\par
\begin{center}
{\footnotesize\begin{tabular}{|p{2.6cm}|p{2.6cm}|c|}
\hline Initial pressure of NO(atm) & Initial pressure of H$_2$(atm) & Rate\\
\hline 0.15 & 0.40 & 0.02\\
\hline 0.075 & 0.40 & 0.005\\
\hline 0.150 & 0.2 & 0.01\\
\hline
\end{tabular}}
\end{center}
Find out the rate law and order of reaction-
\begin{qoptions}
\opt{A}{k[P$_{\mathrm{NO}_2}$]$^2$[P$_{\mathrm{H}_2}$]$^1$, 3}
\opt{B}{k[P$_{\mathrm{NO}_2}$]$^2$[P$_{\mathrm{H}_2}$]$^1$, 2}
\opt{C}{k[P$_{\mathrm{NO}_2}$]$^2$[P$_{\mathrm{H}_2}$]$^1$, 1}
\opt{D}{k[P$_{\mathrm{NO}_2}$]$^2$[P$_{\mathrm{H}_2}$]$^1$, 0}
\end{qoptions}
\qn{13} In dilute aqueous H$_2$SO$_4$, the complex diaquodioxalatoferrate(II) is oxidized by MnO$_4^-$. For this reaction, the ratio of the rate of change of [H$^+$] to the rate of change of [MnO$_4^-$] is:
\begin{qoptions}
\opt{A}{16}
\opt{B}{8}
\opt{C}{4}
\opt{D}{6}
\end{qoptions}
\qn{14} A graph plotted between $\log t_{50\%}$ vs. $\log$ concentration is a straight line. What conclusion can you draw from this graph, (assuming order of reaction is `n')\par
\qfig{p1_chem_s1_q14_1}
\begin{qoptions}
\opt{A}{n = 1 : t$_{1/2}\propto$ a}
\opt{B}{n = 2 : t$_{1/2}\propto$ 1/a}
\opt{C}{n = 1 : t$_{1/2}$ = (0.693/k)}
\opt{D}{None of these}
\end{qoptions}
\qn{15} The forward rate constant for the elementary reversible gaseous reaction C$_2$H$_6\rightleftharpoons$ 2CH$_3$ is $1.57\times10^{-3}$s$^{-1}$ at 100K\par
What is the rate constant for the backward reaction at this temperature if $10^{-4}$ mole of CH$_3$, 10 moles of C$_2$H$_6$ are present in a 10 litre vessel at equilibrium.
\begin{qoptions}
\opt{A}{$1.57\times10^{9}$Lmol$^{-1}$s$^{-1}$}
\opt{B}{$1.57\times10^{10}$Lmol$^{-1}$s$^{-1}$}
\opt{C}{$1.57\times10^{11}$Lmol$^{-1}$s$^{-1}$}
\opt{D}{$1.57\times10^{7}$Lmol$^{-1}$s$^{-1}$}
\end{qoptions}
\qn{16} The rate of reaction depends on concentration as $-\dfrac{dC_A}{dt}=\dfrac{K_1C_A}{1+K_2C_A}$. If C$_A$ being very high, then order of reaction is....
\begin{qoptions}
\opt{A}{1}
\opt{B}{Zero}
\opt{C}{2}
\opt{D}{3}
\end{qoptions}
\qn{17} The conductivity of a solution of AgCl at 298 K is found to be $1.382\times10^{-6}\,\Omega^{-1}$cm$^{-1}$. The ionic conductance of Ag$^+$ and Cl$^-$ at infinite dilution are $61.9\ \Omega^{-1}$cm$^2$mol$^{-1}$ and $76.3\ \Omega^{-1}$cm$^2$mol$^{-1}$, respectively. The solubility of AgCl is
\begin{qoptions}
\opt{A}{$1.4\times10^{-6}$mol L$^{-1}$}
\opt{B}{$1\times10^{-2}$mol L$^{-1}$}
\opt{C}{$1\times10^{-5}$mol L$^{-1}$}
\opt{D}{$1.9\times10^{-5}$mol L$^{-1}$}
\end{qoptions}
\qn{18} Molar conductances of BaCl$_2$, H$_2$SO$_4$ and HCl at infinite dilutions are x$_1$, x$_2$ and x$_3$ respectively. Equivalent conductance of BaSO$_4$ at infinite dilution will be:
\begin{qoptions}
\opt{A}{$\dfrac{[x_1+x_2-x_3]}{2}$}
\opt{B}{$\dfrac{[x_1-x_2-x_3]}{2}$}
\opt{C}{$2\,(x_1+x_2-2x_3)$}
\opt{D}{$\dfrac{[x_1+x_2-2x_3]}{2}$}
\end{qoptions}
\qn{19} When the electric current is passed through a cell having an electrolyte, the positive ions move towards cathode and negative ions towards the anode. If the cathode is pulled out of the solution
\begin{qoptions}
\opt{A}{The positive ions will move towards anode}
\opt{B}{The positive ions will start moving towards the anode while negative ions will stop moving}
\opt{C}{The negative ions will continue to move towards anode while positive ions will stop moving}
\opt{D}{The positive and negative ions will start moving randomly}
\end{qoptions}
\qn{20} If for a reaction in which A(g) converts to B(g) the reaction carried out at const. V \& T results into the following graph.\par
\qfig{p1_chem_s1_q20_1}
\begin{qoptions}
\opt{A}{then the reaction must be A(g) $\to$3B(g) and is a first order reaction.}
\opt{B}{then the reaction must be A(g) $\to$3B(g) and is a second order reaction.}
\opt{C}{then the reaction must be A(g) $\to$3B(g) and is a zero order reaction.}
\opt{D}{then the reaction must be A(g) $\rightleftharpoons$3B(g) and is a first order reaction.}
\end{qoptions}
\sect{SECTION-II}
\qn{1} The decomposition reaction\par
2N$_2$O$_5$ (g) $\xrightarrow{\Delta}$ 2N$_2$O$_4$ (g) + O$_2$ (g)\par
is started in a closed cylinder under isothermal isochoric condition at an initial pressure of 1 atm. After Y$\times10^3$ s, the pressure inside the cylinder is found to be 1.45 atm. If the rate constant of the reaction is $5\times10^{-4}$s$^{-1}$, assuming ideal gas behaviour, the value of Y is
\qn{2} An acid-catalysed hydrolysis of ester is carried out in the presence of HCl. Small but same volume of reaction mixture were taken out at various instant and were titrated with standard NaOH solution.\par
The following data was obtained.\par
\begin{center}
{\footnotesize\begin{tabular}{|p{3.6cm}|c|c|c|}
\hline Time (in min) & 0 & 120 & $\infty$\\
\hline Volume(mL) of NaOH solution needed for titration & 20 & 55 & 60\\
\hline
\end{tabular}}
\end{center}
The value of half-life of reaction is\underline{\hspace{1cm}}min
\qn{3} The magnitude of $E^\circ_{\mathrm{OP}}$ of following reaction is \underline{\hspace{1cm}}volt.\par
Ag(s) + 2NH$_3$ (aq) $\to$ Ag(NH$_3$)$_2^+$ (aq) + e$^-$\par
[Given : K$_f$ (Ag(NH$_3$)$_2^+$) $=10^8$; $E^\circ_{\mathrm{Ag}^+|\mathrm{Ag}}=0.80$V; $2.303\dfrac{RT}{F}=0.06$]
\qn{4} For the cell reaction: Hg$_2$Br$_2$ (s) + 2Ag (s) $\to$ 2Hg ($\ell$) + 2AgBr (s) ,the variation of emf with temperature is shown by the graph given below\par
\qfig{p1_chem_s2_q04_1}
Calculate $\Delta H$ of cell reaction(in kJ) at 300K.
\qn{5} NO$_3^-\to$ NO$_2$(acid medium), E$^\text{o}$=0.790V\par
NO$_3^-\to$ NH$_3$OH$^+$(acid medium) E$^\text{o}$=0.731V.\par
At what pH. the above two will have same E value? Assume the concentration of all other species NH$_3$OH$^+$ except[H$^+$] to be unity.
\qn{6} 3A current was passed through an aqueous solution of an unknown salt of Pd$^{n+}$ for 1Hr. 2.977g of Pd was deposited at cathode. Find n.(Given Atomic mass of Pd=106.4)
\qn{7} Based on the following standard oxidation electrode\par
$E^\circ_{\mathrm{A}^{+2}/\mathrm{A}}=-0.92$V\par
$E^\circ_{\mathrm{B}/\mathrm{B}^{+1}}=+0.18$V\par
$E^\circ_{\mathrm{C}^{+3}/\mathrm{C}}=+0.02$V\par
$E^\circ_{\mathrm{D}/\mathrm{D}^{+2}}=-1.20$V\par
Number of feasible electrochemical cells that can be constructed using metal C as anode at 298K will be
\qn{8} For the reaction 2Fe$^{3+}$(aq) + 2I$^-$(aq) $\to$ 2Fe$^{2+}$(aq) + I$_2$(s)\par
The magnitude of the standard molar free energy change, $\Delta_rG^\circ_m=-$\underline{\hspace{1.5cm}} kJ (Round off to the Nearest Integer).\par
[$E^\circ_{\mathrm{Fe}^{2+}/\mathrm{Fe(s)}}=-0.440$ V; $E^\circ_{\mathrm{Fe}^{3+}/\mathrm{Fe(s)}}=-0.036$ V; $E^\circ_{\mathrm{I}_2/2\mathrm{I}^-}=0.539$ V; $F=96500$ C]
\qn{9} The reaction that occurs in a breath analyser, a device used to determine the alcohol level in a person's blood stream is\par
2K$_2$Cr$_2$O$_7$ + 8H$_2$SO$_4$ + 3C$_2$H$_6$O $\to$ 2Cr$_2$(SO$_4$)$_3$ + 3C$_2$H$_4$O$_2$ + 2K$_2$SO$_4$ + 11H$_2$O\par
If the rate of appearance of Cr$_2$(SO$_4$)$_3$ is 2.67 mol min$^{-1}$ at a particular time, the rate of disappearance of C$_2$H$_6$O at the same time is.........mol min$^{-1}$(Nearest integer)
\qn{10} For the non-equilibrium process,\par
A + B $\to$ Products\par
the rate is first-order with respect to A and second order w.r.t B,\par
when [A] = 1.0M, [B] = 1.0M, rate = $1.0\times10^{-2}$ Ms$^{-1}$,\par
When 75\% each of reactant has been converted into products, rate of reaction is \underline{\hspace{0.8cm}}$\times10^{-4}$Ms$^{-1}$.
\partsec{3}{MATHEMATICS}
\sect{SECTION-I}
\qn{1} $\displaystyle\int x^x(1+\ln x)\,dx=$
\begin{qoptions}
\opt{A}{$x^{\ln x}+c$}
\opt{B}{$x^x+c$}
\opt{C}{$x\ln x+c$}
\opt{D}{None of these}
\end{qoptions}
\qn{2} $\displaystyle\int\dfrac{\cos^2x-\sin^2x}{\cos x+\sin x}\,dx=f(x)$ and if $f(0)=1$, then $f\!\left(\dfrac\pi4\right)=$
\begin{qoptions}
\opt{A}{$\sqrt2$}
\opt{B}{$2+\sqrt2$}
\opt{C}{$2-\sqrt2$}
\opt{D}{0}
\end{qoptions}
\qn{3} If $\displaystyle\int x^2\tan^{-1}x\,dx=\dfrac{x^3}{A}\tan^{-1}x+\dfrac{x^2}{B}$\par
$+\dfrac{\ln(1+x^2)}{C}+k$ then
\begin{qoptions}
\opt{A}{$A+B+C=6$}
\opt{B}{$|ABC|=36$}
\opt{C}{$A+B+C=3$}
\opt{D}{$|ABC|=18$}
\end{qoptions}
\qn{4} If $f(x)=1+3^x\ln3$ and $f(x)$ be its antiderivative which assumes the value of 7 for $x=2$. The value of $x$ for which the curve $f(x)$ cuts the $x$ - axis is
\begin{qoptions}
\opt{A}{0}
\opt{B}{1}
\opt{C}{2}
\opt{D}{3}
\end{qoptions}
\qn{5} $\displaystyle\int\dfrac{\tan^{-1}\sqrt x}{(1+x)\sqrt x}\,dx$
\begin{qoptions}
\opt{A}{$\tan^{-1}x+c$}
\opt{B}{$(\tan^{-1}x)^2+c$}
\opt{C}{$(\tan^{-1}\sqrt x)^2+c$}
\opt{D}{$2(\tan^{-1}\sqrt x)^2+c$}
\end{qoptions}
\qn{6} $\displaystyle\int\dfrac{dx}{x(x^5+1)}=$
\begin{qoptions}
\opt{A}{$5\ln\left|\dfrac{x^5}{x^5+1}\right|+c$}
\opt{B}{$\dfrac15\ln\left|\dfrac{x^5}{x^5+1}\right|+c$}
\opt{C}{$5\ln\left|1+\dfrac1{x^5}\right|+c$}
\opt{D}{$\dfrac15\ln\left|1+\dfrac1{x^5}\right|+c$}
\end{qoptions}
\qn{7} $\displaystyle\int\ln\!\left(x-\sqrt{x^2-1}\right)dx=$\par
$f(x)\ln\!\left(x-\sqrt{x^2-1}\right)+g(x)+c$ where $g(1)=0$, then
\begin{qoptions}
\opt{A}{$f(x)=x^2,\ g(x)=\sqrt{x^2-1}$}
\opt{B}{$f(x)=x,\ g(x)=\sqrt{x^2-1}$}
\opt{C}{$f(x)=x,\ g(x)=\sqrt{1-x^2}$}
\opt{D}{$f(x)=x^2,\ g(x)=\sqrt{1-x^2}$}
\end{qoptions}
\qn{8} If $\displaystyle\int\dfrac{(4\sin x+3\cos x)}{(4\cos x+3\sin x)}dx=ax+b\ln\left|3\sin x+4\cos x\right|+c$,\par
then value of $125(a+b)$ is
\begin{qoptions}
\opt{A}{25}
\opt{B}{45}
\opt{C}{65}
\opt{D}{85}
\end{qoptions}
\qn{9} If $\displaystyle\int\dfrac{dx}{(x+2)(x^2+1)}=a\ln\left(1+x^2\right)+b\tan^{-1}x+0.2\ln|x+2|+c$, then $10a-5b$
\begin{qoptions}
\opt{A}{1}
\opt{B}{$-3$}
\opt{C}{3}
\opt{D}{$-1$}
\end{qoptions}
\qn{10} If $\displaystyle\int\dfrac{(3x+2)\,dx}{4x^2+4x+5}=a\ln\left|4x^2+4x+5\right|+b\tan^{-1}(x+0.5)+c$\par
then $32a+\dfrac ab=$
\begin{qoptions}
\opt{A}{9}
\opt{B}{12}
\opt{C}{15}
\opt{D}{21}
\end{qoptions}
\qn{11} $\displaystyle\int\dfrac{\cos\mathrm{ec}^2x-2023}{(\cos x)^{2023}}\,dx=$
\begin{qoptions}
\opt{A}{$\cot x(\sec x)^{2023}+c$}
\opt{B}{$\tan x(\sec x)^{2023}+c$}
\opt{C}{$-\tan x(\sec x)^{2023}+c$}
\opt{D}{$\dfrac{-\cot x}{(\cos x)^{2023}}+c$}
\end{qoptions}
\qn{12} $\displaystyle\int_1^{\sqrt3}\dfrac{dx}{1+x^2}=\dfrac\pi k$; Then $k=$\underline{\hspace{1cm}}.
\begin{qoptions}
\opt{A}{6}
\opt{B}{3}
\opt{C}{12}
\opt{D}{1}
\end{qoptions}
\qn{13} $\displaystyle\int_{-1}^{1}(x-[2x])\,dx=$ (where $[\,\cdot\,]$ denotes GIF)
\begin{qoptions}
\opt{A}{0}
\opt{B}{1}
\opt{C}{2}
\opt{D}{4}
\end{qoptions}
\qn{14} $\displaystyle\int_2^3\dfrac{(x+2)^2}{2x^2-10x+53}\,dx=$
\begin{qoptions}
\opt{A}{2}
\opt{B}{1}
\opt{C}{0.5}
\opt{D}{2.5}
\end{qoptions}
\qn{15} $\displaystyle\int_{-\pi/2}^{\pi/2}\dfrac{\sin^2x}{1+e^x}\,dx=$
\begin{qoptions}
\opt{A}{$\dfrac\pi8$}
\opt{B}{$\dfrac\pi2$}
\opt{C}{$4\pi$}
\opt{D}{$\dfrac\pi4$}
\end{qoptions}
\qn{16} $\displaystyle\int_0^\pi\dfrac{dx}{2\sin^2x+1}=$
\begin{qoptions}
\opt{A}{$\dfrac\pi4$}
\opt{B}{$\dfrac\pi{\sqrt3}$}
\opt{C}{0}
\opt{D}{$\dfrac\pi{\sqrt2}$}
\end{qoptions}
\qn{17} $\displaystyle\int_{-\pi/4}^{\pi/4}\dfrac{\left(x^9-3x^5+7x^3-x+1\right)}{(1+\cos2x)}\,dx=$
\begin{qoptions}
\opt{A}{6}
\opt{B}{2}
\opt{C}{0}
\opt{D}{1}
\end{qoptions}
\qn{18} $\displaystyle\int_{-2\pi}^{5\pi}\cot^{-1}(\tan x)\,dx=$
\begin{qoptions}
\opt{A}{$\dfrac{7\pi^2}2$}
\opt{B}{$\dfrac{\pi^2}2$}
\opt{C}{0}
\opt{D}{None of these}
\end{qoptions}
\qn{19} Let $I_1=\displaystyle\int_0^1\dfrac{\ln\!\left(\frac1{1+x}\right)}{1+x^2}\,dx$ and $I_2=\displaystyle\int_0^{\pi/2}\ln(\sin x)\,dx$\par
then $\left|\dfrac{I_2}{I_1}\right|=$ ($|\cdot|$ represent modulus)
\begin{qoptions}
\opt{A}{1}
\opt{B}{2}
\opt{C}{3}
\opt{D}{4}
\end{qoptions}
\qn{20} If $\displaystyle\int_0^1\dfrac{\tan^{-1}x}{x}\,dx=k$, then $\displaystyle\int_0^{\pi/2}\dfrac{x}{\sin x}\,dx=$
\begin{qoptions}
\opt{A}{$k$}
\opt{B}{$2k$}
\opt{C}{$\pi\ln2-k$}
\opt{D}{$\pi\ln2+k$}
\end{qoptions}
\sect{SECTION-II}
\qn{1} If $\displaystyle\int\dfrac{2x^{12}+8x^9}{(x^5+x^3+1)^3}\,dx=\dfrac{f(x)}{(x^5+x^3+1)^2}+c$ then the sum of digits of $f(2)$ is \underline{\hspace{1cm}}
\qn{2} $\displaystyle\int\dfrac{11+4\cos2x+\cos4x}{1-\cos4x}\,dx=$\par
$\alpha x+\beta\cot kx+\gamma\,\mathrm{cosec}\,kx+c$\par
Then $\alpha-\beta+k-\gamma=$\underline{\hspace{1cm}}.
\qn{3} $\displaystyle\int\dfrac{(\cos3)\,dx}{\sin(x+5)\cos(x+2)}=\ell n\left|f(x)\right|+c$;\par
$\theta\in N$ then $|f(1)|=a\sin\theta$, then $\theta=$
\qn{4} $\displaystyle\int\dfrac{x^2\,dx}{1+x^6}=f(x)$ If $f(0)=\dfrac\pi3$, then $f(1)=\dfrac{m\pi}n$ where m, n are coprime natural numbers, then m + n = \underline{\hspace{1cm}}.
\qn{5} $\displaystyle\int_{-1}^{1}\log\!\left(x+\sqrt{x^2+1}\right)dx=$\underline{\hspace{1cm}}
\qn{6} $\displaystyle\int_{-\pi}^{\pi}\dfrac{2x(1+\sin x)}{1+\cos^2x}\,dx=\dfrac{\pi^2}k$, then $k=$\underline{\hspace{1cm}}.
\qn{7} $\displaystyle\int_0^2x^{14}\left(x^3-x^2+1\right)^4\left(6x^3-5x^2+3\right)dx=2^p\cdot5^r$\par
then $(p-r)=$\underline{\hspace{1cm}}.
\qn{8} If $g(x)=f(x)+f\!\left(\dfrac1x\right)$, where\par
$f(x)=\displaystyle\int_1^x\dfrac{\ln t}{1+t}\,dt$, then $g(e)$ is\underline{\hspace{1cm}}.
\qn{9} $\displaystyle\int_0^1\dfrac{4\ln(1+x)}{1+x^2}\,dx+\int_0^{\pi/2}\ln(\cos x)\,dx=k$, then $[k]=$\underline{\hspace{1cm}}. (where $[\,\cdot\,]$ is GIF).
\qn{10} Let $f:R\to R$ be a continuous function $f(x)+f(x+1)=2,\ \forall\ x\in R$ \& if\par
$I_1=\displaystyle\int_0^8f(x)\,dx$ and $I_2=\displaystyle\int_{-1}^3f(x)\,dx$, then $I_1+2I_2=$\underline{\hspace{1cm}}.
\end{multicols}
\end{document}
